\documentclass{amsart}
\usepackage{amsmath,amssymb,amsthm,hyperref}
\newtheorem{theorem}{Theorem}
\newtheorem{lemma}{Lemma}
\begin{document}

\section*{Solution}

\subsection*{Statement}
We seek all quadruples of perfect squares $A=x^2$, $B=y^2$, $C=z^2$, $D=t^2$,
with $x,y,z,t$ positive integers (or, more generally, positive rationals),
such that the four successive products
\[
P_1 = A\cdot A,\quad P_2 = A\cdot A\cdot B,\quad P_3 = A\cdot A\cdot B\cdot C,\quad
P_4 = A\cdot A\cdot B\cdot C\cdot D
\]
satisfy $(P_1,P_2,P_3,P_4)=(A,B,C,D)$. In terms of $x,y,z,t$ this is the system
\begin{align}
x^4 &= x^2, \tag{E1}\\
x^4 y^2 &= y^2, \tag{E2}\\
x^4 y^2 z^2 &= z^2, \tag{E3}\\
x^4 y^2 z^2 t^2 &= t^2. \tag{E4}
\end{align}

\subsection*{Result}

\begin{theorem}
Let $x,y,z,t$ be positive real numbers (in particular positive integers or positive
rationals). Then the system \textup{(E1)--(E4)} holds if and only if
\[
x=y=z=1,
\]
while $t$ is \emph{arbitrary}. Equivalently, the complete set of admissible quadruples
of squares is
\[
(A,B,C,D)=(1,\,1,\,1,\,t^2),\qquad t>0 .
\]
In the integer setting the solutions are exactly $(A,B,C,D)=(1,1,1,n^2)$ with
$n\in\mathbb{Z}_{\ge 1}$; in the positive-rational setting they are
$(1,1,1,r^2)$ with $r\in\mathbb{Q}_{>0}$. In particular the only solution in which all
four squares are forced to a single value is the trivial one $(1,1,1,1)$, and every
non-trivial solution is obtained from it by leaving the last square $D=t^2$ free.
\end{theorem}

\begin{proof}
Throughout we use that all the quantities $x,y,z,t$ are \emph{positive}; hence
$x^2,y^2,z^2,t^2$ are positive and may be cancelled from both sides of an equation.

\medskip
\noindent\textbf{Step 1: (E1) forces $x=1$.}
Equation (E1) reads $x^4=x^2$, i.e.
\[
x^2\,(x^2-1)=0 .
\]
Since $x>0$ we have $x^2>0$, so $x^2-1=0$, giving $x^2=1$ and therefore $x=1$
(the positive root). Consequently
\[
x^2=1,\qquad x^4=1 ,
\]
so $A=x^2=1$ and $P_1=x^4=1=A$, as required.

\medskip
\noindent\textbf{Step 2: substitute $x=1$ and reduce.}
With $x=1$ we have $x^4=1$, and the remaining equations simplify to
\begin{align}
y^2 &= y^2, \tag{E2$'$}\\
y^2 z^2 &= z^2, \tag{E3$'$}\\
y^2 z^2 t^2 &= t^2. \tag{E4$'$}
\end{align}
Equation (E2$'$) is an identity and imposes no constraint by itself.

\medskip
\noindent\textbf{Step 3: (E3$'$) forces $y=1$.}
Dividing (E3$'$) by $z^2>0$ gives $y^2=1$, hence $y=1$ (positive root).
Thus $B=y^2=1$.

\medskip
\noindent\textbf{Step 4: (E4$'$) forces $z=1$.}
With $y=1$, equation (E4$'$) becomes $z^2 t^2=t^2$. Dividing by $t^2>0$ gives
$z^2=1$, hence $z=1$ (positive root). Thus $C=z^2=1$.

\medskip
\noindent\textbf{Step 5: $t$ is unconstrained.}
After $x=y=z=1$, all four equations are satisfied for \emph{every} positive $t$:
indeed $x^4=1$, so
\[
P_1=1=A,\quad P_2=1\cdot 1=1=B,\quad P_3=1\cdot1\cdot1=1=C,\quad
P_4=1\cdot1\cdot1\cdot t^2=t^2=D .
\]
No equation involves $t$ on only one side, so $t>0$ is free; equivalently the last
square $D=t^2$ may be any positive square.

\medskip
\noindent\textbf{Conclusion.}
Steps 1--4 show that any solution must have $x=y=z=1$, and Step 5 shows that
$x=y=z=1$ with arbitrary $t>0$ is indeed a solution. Hence the solution set is exactly
\[
\{(x,y,z,t):x=y=z=1,\ t>0\},
\]
that is $(A,B,C,D)=(1,1,1,t^2)$ with $t>0$. Restricting $t$ to positive integers
(resp.\ positive rationals) yields the integer (resp.\ rational) solutions stated in
the theorem.
\end{proof}

\subsection*{Remarks}
\begin{itemize}
\item The asymmetry of the answer (the parameters $x,y,z$ are pinned to $1$ while $t$
remains free) reflects the structure of the system: each successive product equation
$P_{k}=$ (the $k$-th square) cancels one more variable, but the \emph{last} square
$D=t^2$ appears on \emph{both} sides of the final equation~(E4) and therefore is not
constrained once the earlier variables equal $1$.
\item In particular there is \emph{no} non-trivial quadruple with all four entries
distinct and $\neq 1$: the only freedom is in $D$.
\end{itemize}

\end{document}
